\chapter*{\centering Abstract}
\addcontentsline{toc}{chapter}{Abstract}


Single atoms coupled to optical cavities provide a powerful platform for photonic quantum information processing.
Minimizing the influence of external factors like magnetic field fluctuations is a key goal in this endeavor.
In this work, I investigate a novel gate protocol that employs a single $^{87}Rb$ atom prepared in a magnetic-field-insensitive clock qubit and coupled to a high-finesse birefringent Fabry–Perot cavity.
The cavity’s high birefringence creates two well-separated polarization eigenmodes, enabling polarization-selective reflectivity that depends on the atomic qubit state.
This allows for the implementation of a controlled phase (CPHASE) gate between the atomic clock states and our photonic qubits.
By rotating the basis of the photonic qubit, one is then able to realize an atom-photon controlled-NOT (CNOT) gate.
\\

\textbf{Show some results here and talk about how well it works and also how one might be able to use the gate}
\\


I present simulations based on input–output theory and master-equation modeling that quantify the conditional reflection amplitudes, gate truth table, and resulting fidelities under realistic experimental imperfections such as finite mode matching, cavity loss channels, and multiphoton contributions.
Our results indicate that atom-state-dependent phase shifts on the photonic qubit are achievable in the current system, providing a viable path toward a robust, high-fidelity, cavity-assisted atom–photon quantum gate.

